
\label{secannexes}
\section{Annexes}

\begin{frame}
    \frametitle{Annexes}
    \tableofcontents[
        currentsection,
        sections=\value{section},
        subsubsectionstyle=hide
    ]
\end{frame}
\addtocounter{framenumber}{-1}

\subsection{\textit{Rightmost derivation}}

\begin{frame}
    \frametitle{\textit{Rightmost derivation}}

    \begin{minipage}{0.5\textwidth}
        \[
    \begin{aligned}
        \axiomcolour{E} & \to \axiomcolour{E}+\axiomcolour{E} \\
        \axiomcolour{E} & \to \axiomcolour{E}\times\axiomcolour{E} \\
        \axiomcolour{E} & \to \text{int}
    \end{aligned}
\]

    \end{minipage}%
    \begin{minipage}{0.5\textwidth}
        \[
            \uncover<3->{ \rlap{$\overbrace{\phantom{1 + 2}}^{\mathclap{\text{\textit{leftmost}}}}$}}
            1 + \onslide<3-> \underbrace{ \onslide<1->
                2 \times 3
            \onslide<3-> }_{\mathclap{\text{\textit{rightmost}}}} \onslide<1->
        \]
    \end{minipage}
    \vspace{1em}

    \uncover<2->{
        \begin{minipage}{0.5\textwidth}
            \begin{center}
                \begin{forest}
    [+
        [1]
        [$\times$
            [2]
            [3]
        ]
    ]
\end{forest}


                \textit{Leftmost}

                \uncover<4->{Analyse descendante}

                \uncover<5->{LL(1)}
            \end{center}
        \end{minipage}%
        \begin{minipage}{0.5\textwidth}
            \begin{center}
                \begin{forest}
    [$\times$
        [+
            [1]
            [2]
        ]
        [3]
    ]
\end{forest}


                \textit{Rightmost}

                \uncover<4->{Analyse ascendante}

                \uncover<5->{LR(0), LR(1), LALR(1), \dots}
            \end{center}
        \end{minipage}
    }
\end{frame}
\addtocounter{framenumber}{-1}

\subsection{Limitations de LR(0)~: exemples de conflits}

\begin{frame}
    \frametitle{Limitations de LR(0)~: exemples de conflits}

    \[
    \begin{aligned}
        \axiomcolour{IF\_EXPR} & \to \text{si bool}~\ntcolour{EXPR} \\
        \axiomcolour{IF\_EXPR} & \to \text{si bool}~\ntcolour{EXPR}~\text{sinon}~\ntcolour{EXPR} \\
        \ntcolour{EXPR}        & \to \text{instruction}
    \end{aligned}
\]

    \pause
    \begin{tikzpicture}[
        every node/.style={rectangle, font=\footnotesize, inner sep=3pt},
        every state/.style={rectangle, font=\footnotesize, inner sep=1pt},
        every loop/.style={min distance=1mm, looseness=2},
        auto,
        node distance=7mm and 7mm
    ]

    \node (q1) [state, initial, initial text={}] {
        $\begin{aligned}
            \axiomcolour{I} & \to \pointeur \text{si bool}~\ntcolour{E} \\
            \axiomcolour{I} & \to \pointeur \text{si bool}~\ntcolour{E}~\text{sinon}~\ntcolour{E}
        \end{aligned}$
    };

    \node (q2) [state, below=of q1] {
        $\begin{aligned}
            \axiomcolour{I} & \to \text{si \pointeur bool}~\ntcolour{E} \\
            \axiomcolour{I} & \to \text{si \pointeur bool}~\ntcolour{E}~\text{sinon}~\ntcolour{E}
        \end{aligned}$
    };

    \node (q3) [state, below=of q2] {
        $\begin{aligned}
            \axiomcolour{I} & \to \text{si bool} \pointeur \ntcolour{E} \\
            \axiomcolour{I} & \to \text{si bool} \pointeur \ntcolour{E}~\text{sinon}~\ntcolour{E} \\
            \ntcolour{E}    & \to \pointeur \text{instr}
        \end{aligned}$
    };

    \node (q4) [state, right=of q3, xshift=30.5mm] {
        $\begin{aligned}
            \axiomcolour{I} & \to \text{si bool}~\ntcolour{E} \pointeur \\
            \axiomcolour{I} & \to \text{si bool}~\ntcolour{E} \pointeur \text{sinon}~\ntcolour{E}
        \end{aligned}$
    };

    \uncover<3-> {
        \node (boxq4) [draw, inner sep=1mm, red, thick, fit=(q4)] {};
    }

    \node (q5) [state, right=of q2, xshift=30mm] {
        $\begin{aligned}
            \axiomcolour{I} & \to \text{si bool}~\ntcolour{E}~\text{sinon} \pointeur \ntcolour{E} \\
            \ntcolour{E}    & \to \pointeur \text{instr}
        \end{aligned}$
    };

    \node (q6) [state, above=of q5] {
        $\begin{aligned}
            \axiomcolour{I} & \to \text{si bool}~\ntcolour{E}~\text{sinon}~\ntcolour{E} \pointeur
        \end{aligned}$
    };

    \node (q7) [state, right=of q2, xshift=4mm] {
        $\ntcolour{E} \to \text{instr} \pointeur$
    };

    \path[->, every node/.style={font=\footnotesize}]
        (q1)  edge node              {si}           (q2)
        (q2)  edge node              {bool}         (q3)
        (q3)  edge node              {\ntcolour{E}} (q4)
        (q4)  edge node              {sinon}        (q5)
        (q5)  edge node              {\ntcolour{E}} (q6)
        (q3)  edge node [auto=right] {instr}        (q7)
        (q5)  edge node [auto=right] {instr}        (q7);

\end{tikzpicture}

\end{frame}
\addtocounter{framenumber}{-1}

\begin{frame}
    \frametitle{Limitations de LR(0)~: exemples de conflits}

    \begin{center}
        \begin{tikzpicture}[
        every node/.style={rectangle, inner sep=3pt},
    ]

    \node (q4) [draw] {
        $\begin{aligned}
            \underline{I} & \to \text{si bool E \pointeur} \\
            \underline{I} & \to \text{si bool E \pointeur sinon E}
        \end{aligned}$
    };

\end{tikzpicture}

    \end{center}

    Que faire~?
    \pause
    \begin{itemize}
        \item Réduire $\axiomcolour{I} \to \text{\textit{si bool \ntcolour{E}}}$~?
        \pause
        \item Essayer de lire $sinon$~?
    \end{itemize}

    \pause
    C’est un conflit lire/réduire.
\end{frame}
\addtocounter{framenumber}{-1}

\begin{frame}
    \frametitle{Limitations de LR(0)~: différence avec LR(1)}

    \begin{center}
        \begin{tikzpicture}[
        every node/.style={rectangle, inner sep=3pt},
    ]

    \node (q4) [draw] {
        $\begin{aligned}
            \axiomcolour{I} & \to \text{si bool}~\ntcolour{E} \pointeur ~\wr \{\lozenge\} \\
            \axiomcolour{I} & \to \text{si bool}~\ntcolour{E} \pointeur \text{sinon}~\ntcolour{E} ~\wr \{\text{\lozenge}\}
        \end{aligned}$
    };

\end{tikzpicture}

    \end{center}

    \pause
    On regarde le lexème suivant~:
    \pause
    \begin{itemize}
        \item Si c’est la fin du fichier (\lozenge), réduire
            $\axiomcolour{I} \to \text{\textit{si bool \ntcolour{E}}}$~;
        \pause
        \item Si c’est $sinon$, lire.
    \end{itemize}
\end{frame}
\addtocounter{framenumber}{-1}

\begin{frame}
    \frametitle{Limitations de LR(0)~: conflits réduire/réduire}

    \begin{center}
        \begin{tikzpicture}[
        every node/.style={rectangle, inner sep=3pt},
    ]

    \node (q4) [draw] {
        $\begin{aligned}
            \ntcolour{A} & \to \dots \pointeur \\
            \ntcolour{B} & \to \dots \pointeur
        \end{aligned}$
    };

\end{tikzpicture}

    \end{center}

    Deux choix~:
    \pause
    \begin{itemize}
        \item Réduire $\ntcolour{A} \to \dots$~;
        \pause
        \item Réduire $\ntcolour{B} \to \dots$.
    \end{itemize}

    \pause
    C’est un conflit réduire/réduire.
\end{frame}
\addtocounter{framenumber}{-1}

\begin{frame}
    \frametitle{Limitations de LR(0)~: conflits réduire/réduire}

    \begin{center}
        \begin{tikzpicture}[
        every node/.style={rectangle, inner sep=3pt},
    ]

    \node (q4) [draw] {
        $\begin{aligned}
            \ntcolour{A} & \to \dots \pointeur ~\wr \{a, b, \lozenge\} \\
            \ntcolour{B} & \to \dots \pointeur ~\wr \{c, d\}
        \end{aligned}$
    };

\end{tikzpicture}

    \end{center}

    On regarde le lexème suivant~:
    \pause
    \begin{itemize}
        \item Si c’est $a$, $b$ ou la fin du fichier (\lozenge), réduire
            $\ntcolour{A} \to \dots$~;
        \pause
        \item Si c’est $c$ ou $d$, réduire $\ntcolour{B} \to \dots$.
    \end{itemize}
\end{frame}
\addtocounter{framenumber}{-1}

\subsection{Automate LR(1)}

\begin{frame}
    \frametitle{Automate LR(1)}

    \[
    \begin{aligned}
        \axiomcolour{IF\_EXPR} & \to \text{si bool}~\ntcolour{EXPR} \\
        \axiomcolour{IF\_EXPR} & \to \text{si bool}~\ntcolour{EXPR}~\text{sinon}~\ntcolour{EXPR} \\
        \ntcolour{EXPR}        & \to \text{instruction}
    \end{aligned}
\]

    \pause
    \maxsizebox{\textwidth}{0.8\textheight}{
        \begin{tikzpicture}[
        every node/.style={rectangle, font=\footnotesize, inner sep=3pt},
        every state/.style={rectangle, font=\footnotesize, inner sep=1pt},
        every loop/.style={min distance=1mm, looseness=2},
        auto,
        node distance=7mm and 7mm
    ]

    \node (q1) [state, initial, initial text={}] {
        $\begin{aligned}
            \underline{I} & \to \pointeur \text{si bool E} ~\wr \{\lozenge\} \\
            \underline{I} & \to \pointeur \text{si bool E sinon E} ~\wr \{\lozenge\}
        \end{aligned}$
    };

    \node (q2) [state, below=of q1] {
        $\begin{aligned}
            \underline{I} & \to \text{si \pointeur bool E} ~\wr \{\lozenge\} \\
            \underline{I} & \to \text{si \pointeur bool E sinon E} ~\wr \{\lozenge\}
        \end{aligned}$
    };

    \node (q3) [state, below=of q2] {
        $\begin{aligned}
            \underline{I} & \to \text{si bool \pointeur E} ~\wr \{\lozenge\} \\
            \underline{I} & \to \text{si bool \pointeur E sinon E} ~\wr \{\lozenge\} \\
            \text{E}      & \to \pointeur \text{instr} ~\wr \{\lozenge, sinon\}
        \end{aligned}$
    };

    \node (q4) [state, right=of q3, xshift=40mm] {
        $\begin{aligned}
            \underline{I} & \to \text{si bool E \pointeur} ~\wr \{\lozenge\} \\
            \underline{I} & \to \text{si bool E \pointeur sinon E} ~\wr \{\lozenge\}
        \end{aligned}$
    };

    \node (q5) [state, right=of q2, xshift=40mm] {
        $\begin{aligned}
            \underline{I} & \to \text{si bool E sinon \pointeur E} ~\wr \{\lozenge\} \\
            \text{E}      & \to \pointeur \text{instr} ~\wr \{\lozenge\}
        \end{aligned}$
    };

    \node (q6) [state, above=of q5, yshift=0.5mm] {
        $\begin{aligned}
            \underline{I} & \to \text{si bool E sinon E \pointeur} ~\wr \{\lozenge\}
        \end{aligned}$
    };

    \node (q7) [state, right=of q2] {
        $\text{E} \to \text{instr} \pointeur ~\wr \{\lozenge, sinon\}$
    };

    \node (q8) [state, left=of q6, xshift=-3mm] {
        $\text{E} \to \text{instr} \pointeur ~\wr \{\lozenge\}$
    };

    \path[->, every node/.style={font=\footnotesize}]
        (q1)  edge              node                           {si}    (q2)
        (q2)  edge              node                           {bool}  (q3)
        (q3)  edge              node                           {E}     (q4)
        (q4)  edge              node                           {sinon} (q5)
        (q5)  edge              node                           {E}     (q6)
        (q3)  edge              node [auto=right]              {instr} (q7)
        (q5)  edge              node [yshift=1mm]              {instr} (q8);

\end{tikzpicture}

    }
\end{frame}
\addtocounter{framenumber}{-1}

\subsection{Nécessité de la récursivité terminale}

\begin{frame}
    \frametitle{Nécessité de la récursivité terminale}

    \begin{center}
        \begin{minipage}{0.5\textwidth}
            \only<1>{
                \inputminted{ocaml}{exemples/0intro/somme_liste_rec_court.ml}
            }%
            \only<2>{
                \inputminted{ocaml}{exemples/3boucles/probleme.ml}
            }
        \end{minipage}
    \end{center}
\end{frame}
\addtocounter{framenumber}{-1}

\subsection{Exemple de fonction sans accumulateur}

\begin{frame}
    \frametitle{\mintinline{ocaml}{palindrome}}

    \begin{adjustwidth}{-1.5 em}{-1.5em}
    \only<1>{\inputminted{ocaml}{exemples/5annexes/palindrome_rec.ml}}
    \only<2>{\inputminted[fontsize=\scriptsize]{ocaml}{exemples/5annexes/palindrome_loop.ml}}
    \end{adjustwidth}
\end{frame}
\addtocounter{framenumber}{-1}
